\section{Data}
\label{sec:data}

\subsection{Source and Sample}

We use one-minute OHLCV bars (open, high, low, close, volume) from Yahoo Finance for
\RELIANCEDays{} NSE trading days, 30 September to 28 October 2026. Yahoo keeps one-minute
history for about 30 days, so this is the full window available. Two cleaning steps matter:
\begin{enumerate}
\item \textbf{Session window.} Yahoo drops the 15:16--15:27 bars and then prints a single
  catch-up bar at 15:28. Keeping that bar would create a fake volatility spike at the close.
  We use the clean window 09:15--15:14, which gives exactly $T=\Horizon$ one-minute bars per day.
\item \textbf{Minutes with no trade.} A missing bar inside the window means nobody traded that
  minute. We carry the last close forward (the mid did not move), leave the high and low empty
  (no quote could have been filled) and set volume to zero. Nothing is interpolated.
\end{enumerate}
The first \CalibDays{} days are the \emph{calibration window}; the remaining
\RELIANCETestDays{} days are held out as the \emph{test window} (Figure~\ref{fig:prices}).

\subsection{Choosing a Liquid and an Illiquid Stock}

The model's predictions depend on liquidity: a thin book means fewer fills per unit of price risk,
so inventory control should matter more. We therefore study a contrasting pair. RELIANCE is the
most heavily traded NSE large-cap. For the illiquid stock we screened twelve small-caps by
Amihud's~\cite{amihud2002} illiquidity ratio and by the share of minutes with no trade, and chose
KSCL (Kaveri Seed). It is among the most illiquid names screened, it still trades in most
minutes, and its price level is close enough to RELIANCE's that spreads in rupees are comparable.
Table~\ref{tab:descriptive} and Figure~\ref{fig:liquidity} show how far apart the two are: KSCL's
Amihud ratio is about 300 times RELIANCE's, and \KSCLNoTrade\% of its minutes have no trade.

\begin{table}[!htbp]
\centering
\caption{Descriptive statistics, all \RELIANCEDays{} days, 09:15--15:14.}
\label{tab:descriptive}
\small
\begin{tabular}{lrr}
\toprule
Statistic & RELIANCE & KSCL \\
\midrule
Trading days & 19 & 19 \\
Mean price (INR) & 1,189.3 & 712.4 \\
Mean daily traded value (INR crore) & 1,640.27 & 9.19 \\
Minutes with no trade (\%) & 0.2\% & 8.8\% \\
$\sigma$, one-minute price change (INR) & 0.823 & 1.001 \\
Mean one-minute high--low range (INR) & 1.020 & 0.655 \\
Amihud ILLIQ ($|R|$ per INR crore) & 7.18e-06 & 2.35e-03 \\
Roll spread (INR) & 0.373 & 0.686 \\
\bottomrule
\end{tabular}

\end{table}

\begin{figure}[!htbp]
\centering
\includegraphics[width=\textwidth]{data_price_paths}
\caption{One-minute closing prices, all sessions placed end to end (thin vertical lines separate
days). The shaded part is the calibration window; the model is tested only on the unshaded days.}
\label{fig:prices}
\end{figure}

\begin{figure}[!htbp]
\centering
\includegraphics[width=\textwidth]{data_liquidity}
\caption{Daily liquidity. (a) Amihud~\cite{amihud2002} illiquidity, $|R_d|$ per crore of traded
value (log scale): KSCL is two to three orders of magnitude less liquid. (b) Roll~\cite{roll1984}
spread, $2\sqrt{-\mathrm{Cov}(\Delta p_t,\Delta p_{t-1})}$; the measure is undefined (gaps) on days
when the covariance is positive, which happens often for RELIANCE. (c) Share of minutes with no
trade.}
\label{fig:liquidity}
\end{figure}

\subsection{The Trading Day Is Not Uniform}

Figure~\ref{fig:volume} shows the familiar intraday U-shape~\cite{admati1988}: trading is heaviest
at the open and in the last half hour. KSCL's curve is noisier and has bursts in the middle of the
day, typical of a stock where a few large orders dominate volume. This pattern matters for the
model, which assumes order flow and volatility are constant through the session (A5).

\begin{figure}[!htbp]
\centering
\includegraphics[width=\textwidth]{data_intraday_volume}
\caption{Average share of daily volume traded in each minute (thin: raw; thick: 9-minute moving
average). The dashed line is a perfectly uniform day.}
\label{fig:volume}
\end{figure}
